\documentclass{report}
\usepackage{graphicx} % Required for inserting images
\usepackage[a4paper, left=2cm, right=2cm, top=3cm, bottom=3cm]{geometry}
\usepackage{amsmath} %usa questo per scrivere le formule matematiche
\usepackage{parskip} %Questo pack elimina l'indentazione fastidiosa 
\usepackage{amsthm} % Questo pack permette di scrivere in modo carino i teoremi,
\usepackage{amssymb}
\usepackage{xcolor}    % Per i colori
\usepackage{framed}    % Per le cornici
\usepackage{mdframed}  % Cornici più avanzate, definizioni, proposizioni ecc.
\usepackage{hyperref}  % Per inserire i link

%Quelli che seguono sono i risultati dell'uso del pack amsthm
\newtheorem{theo}{Theorem}
\newtheorem{prop}{Proposition}
\newtheorem{crlr}{Corollario}[theo] % Numerato in base ai teoremi
\newtheorem{defi}{Definition}
\newtheorem{note}{Note}
\newtheorem{xpl}{Example}
\newenvironment{pop}{}{}

%dati iniziali
\title{Language Generation in the Limit}
\author{Fabio Colletti}
\date{October 2026}

\begin{document}

\maketitle

\chapter{Introduction}
La tesi si basa sulla descrizione e discussione dei risultati dell'articolo \href{https://arxiv.org/pdf/2404.06757}{Language Generation in the Limit} \\
di Jon Kleinberg e Sendhil Mullainathan.
\chapter{Notazione e Definizioni}
\begin{pop}
  In questa sezione vediamo quali sono le notazioni utilizzate nel corso dell'elaborato, e le definizioni non banali.
\end{pop}
\end{document}
